\section{Neuron Architecture}
\label{sec:arch}

Neuron is middleware between an agent framework and host/container services.
The architecture makes one research contribution---a coupled resource and
permission decision path---and separates it from supporting safety
invariants and prototype engineering. Figure-level breadth is deliberately
avoided: quota trees, sandboxes, model routing, and audit storage are not
presented as independent innovations.

\subsection{Control-Plane Boundary}

Tenant agents issue typed runtime requests containing a tenant identifier,
agent identifier, turn/trace identifier, spawn depth, action, and payload.
The dispatcher first applies resource admission, then evaluates permission,
then dispatches an admitted operation to a tool or sandbox. The order is
security-relevant: rejected tenant work should not occupy permission workers,
and an action cannot execute before a permission verdict exists.

Neuron distinguishes two execution domains:

\begin{description}
  \item[Tenant domain.] Agent turns, model callbacks, and tool preparation
  may be concurrent and are assumed adversarial. Their executor and queues
  are capacity limited.
  \item[Policy domain.] Permission evaluation uses reserved workers and a
  bounded queue. It consumes immutable request context plus authenticated
  resource state. The core runtime prototype uses an independent JVM
  executor; the fault-injection artifact also realizes this boundary as a
  separately started Java process behind a loopback HTTP service.
\end{description}

Executor reservation is not coupling. It is a prerequisite for a dependable
policy plane. Coupling occurs only when resource state becomes an input to
the permission function.

\subsection{Core Mechanism: Resource-Informed Permission}
\label{sec:arch-coupling}

Let $d$ be spawn depth, $a$ the requested action, $D$ the normal destructive
action boundary, $r$ the resource-layer rejection count in the configured
window, and $\tau$ the pressure threshold. Let
$\mathit{destructive}(a)$ be the runtime's action classification. The static
policy is

\begin{equation}
P_0(d,a)=
\begin{cases}
\textsc{deny-depth}, & d \ge D \land \mathit{destructive}(a),\\
\textsc{ask-with-context}, & \text{otherwise}.
\end{cases}
\end{equation}

Neuron adds a pressure-conditioned clause:

\begin{equation}
P_c(d,a,r)=
\begin{cases}
\textsc{deny-depth}, & \begin{array}{l}
  \mathit{destructive}(a)\land d\ge D\\[-1pt]
  {}\lor\mathit{destructive}(a)\land r>\tau\land d\ge D-1
\end{array},\\
\textsc{ask-with-context}, & \text{otherwise}.
\end{cases}
\end{equation}

The rule tightens by one level; it never relaxes a static denial. In the
prototype $D=2$. Thus a depth-2 destructive request is denied with or without
coupling. A depth-1 request is the discriminating probe: it is
\textsc{ask-with-context} when $r\le\tau$ and \textsc{deny-depth} when
$r>\tau$. Non-destructive actions are unaffected by this rule.

The Java runtime's \texttt{EscalationPolicy} implements this function. The
resource signal is produced by the runtime's
\texttt{RateLimitSyscallFilter}, a tenant-scoped token bucket. The current
signal is deliberately simple and has limitations: a cumulative count
requires an explicit reset/window policy, may not transfer across hardware,
and may be diluted by Sybil tenants. The evaluation therefore freezes the
counter semantics and treats $\tau$ as a deployment parameter.

\subsection{Policy Availability and Timeout Semantics}

A caller submits the immutable permission task to the policy executor and
waits for a configured deadline. There are three possible outcomes:

\begin{enumerate}
  \item the queue accepts the task and the policy returns before the
  deadline;
  \item the queue rejects the task because it is full; or
  \item the queue accepts the task but it does not finish before the
  deadline.
\end{enumerate}

Outcomes 2 and 3 are distinct telemetry events but require the same explicit
deployment choice. A fail-open wrapper returns a permissive result; a
fail-closed wrapper denies. Neuron's recommended configuration is reserved
capacity plus fail-closed timeout handling. The fail-open configuration is
retained only as an experimental baseline because it exposes the security
consequence under study.

Reservation has two parameters: worker count and queue capacity. Too little
capacity denies legitimate bursts; an unbounded queue converts overload into
unbounded latency and memory consumption. The prototype consequently uses a
bounded queue and records its depth with every observation. The process-
separated fault-injection path provides a stronger CPU/memory boundary but
introduces RPC and partition semantics; the current experiment keeps the
service on loopback, so WAN-scale behavior remains outside the claim.

\subsection{Causal Context and Audit Correlation}

Admission and permission events carry the same trace identifier. A resource
record includes tenant, action, bucket outcome, rejection count, and time. A
permission record includes depth, action class, pressure value consumed,
verdict, executor admission result, deadline outcome, and policy latency.
The pair answers a question that separate logs cannot answer reliably:
which resource state actually informed a particular verdict?

Correlation is supporting infrastructure, not enforcement. If the resource
counter is unauthenticated or a trace identifier can be replaced by tenant
code, the joint rule becomes attacker-controlled. Neuron creates both inside
the runtime boundary and propagates immutable copies to child tasks. Audit
failure does not silently convert a denial to an allow; a deployed system
must choose whether audit backpressure itself is fail-closed or buffered in
reserved storage.

\subsection{Supporting Invariant: Privilege Non-Increase}

Dynamic spawn introduces a second way to cross the policy boundary. Let
$\Pi(x)$ denote the effective capabilities of agent $x$. For every spawn edge
$p\rightarrow c$, Neuron requires

\begin{equation}
\Pi(c) \subseteq \Pi(p).
\end{equation}

The spawn path captures the parent's effective profile before asynchronous
dispatch, derives the child profile by intersection with the requested
profile, and propagates tenant, trace, and depth in the same immutable spawn
context. Missing parent context is not interpreted as a default permissive
profile. Cleanup occurs in a \texttt{finally} block so pooled threads cannot
leak identity to later tasks.

The invariant supports coupled governance because the depth and tenant used
by $P_c$ must describe the actual caller. It is nevertheless orthogonal to
the starvation claim: a system can preserve privilege inheritance while its
policy executor is unavailable. Neuron tests the invariant separately and
does not fold those tests into the response-surface sample count.

\subsection{Engineering Components and Trust Boundary}

The broader prototype includes hierarchical quotas and sandbox providers.
Quotas supply admission decisions; sandboxes contain tools after permission
has been granted. Neither changes the equations above. For clarity:

\begin{itemize}
  \item token quotas and token buckets are resource-admission mechanisms;
  \item a WASM or container provider is an execution-isolation mechanism;
  \item identity propagation, privilege intersection, and correlated audit
  are supporting security mechanisms; and
  \item reserved policy capacity plus $P_c$ is the evaluated coupled
  governance mechanism.
\end{itemize}

\label{sec:arch-sandbox}
The container/JVM boundary remains trusted. This paper neither evaluates
container escape nor claims that an in-process sandbox withstands arbitrary
native code. Newly registered native tools must execute outside the policy
JVM under an independently configured sandbox.

\subsection{Request Flow}

For a destructive tool request, the complete flow is:

\begin{enumerate}
  \item validate the runtime-authenticated tenant, trace, and spawn context;
  \item apply tenant-scoped resource admission and update the pressure state;
  \item submit a permission task to reserved capacity;
  \item compute $P_c$ from depth, action class, and the captured pressure
  value;
  \item apply fail-closed handling if the task cannot be admitted or misses
  its deadline;
  \item append correlated resource and permission records; and
  \item dispatch only an admitted operation to the configured sandbox/tool.
\end{enumerate}

The evaluation manipulates steps 2--5 while keeping the runtime admission
and policy functions fixed. This is what permits attribution among load,
isolation, fallback semantics, and coupling.
